\section{Extension: the laminar column at low carriers}\label{app:lanmm}
The two-band laminar column (LaNMM) adds a resonant intermediary
between the field and the slow generator. These sweeps use $25$--$80$~Hz carriers
and polarizations up to $3$~mV on $P_2$, exceeding the sigmoid scale
$1/\rho\approx1.8$~mV. They describe nonlinear frequency selection below TI's
kHz range. The weak-field square law is tested separately in JR at
$0.05$--$0.2$~mV, one to two orders of magnitude below that scale
(Fig.~\ref{fig:opp}b).

\subsection{Model and measurements}\label{sec:lanmm-methods}
The laminar extension (LaNMM)
\cite{ruffini2025cfc,sanchezTodo2023} adds coupled alpha and gamma circuits: deep
pyramidal $P_1$, spiny-stellate (SS), and somatostatin (SST) populations, and a
superficial pyramidal $P_2$/PV circuit.
We simulate this laminar neural mass model (LaNMM). We integrate the full 28-variable column
\cite{ruffini2025cfc,sanchezTodo2023}, with the field entering as in the single
JR column, as a membrane polarization added inside the sigmoid argument of the target
population. Its output rate $\varphi_X$ (s$^{-1}$) is
\begin{equation}
\varphi_X(t)=\Sig\!\big(v_X(t)+s(t)\big),\qquad X\in\{P_1,P_2\},
\label{eq:lanmm_drive}
\end{equation}
with $s(t)$ the AM or two-tone field of \S\ref{sec:coupling} (amplitude $\varepsilon$ in
mV, beat $\Omega=2\pi\Delta f$, carrier $\omega_c=2\pi f_c$) and all three external drives
held flat at their baselines (Fig.~\ref{fig:lanmm_setup}). A perturbation of an external
input \emph{rate} would require clipping at zero
once its excursion exceeds the baseline, and a hard clip is a half-wave rectifier---a
diode detector placed upstream of the sigmoid, which would inject a line at $\Delta f$
into the drive before the population nonlinearity received it. A signed membrane
polarization carries no such constraint and none is applied. Because the LaNMM and JR
sigmoids share the voltage scale $1/\rho\approx1.8$~mV, $\varepsilon$ is directly
comparable between the two models.

Integration is a vectorized fixed-step RK4
($\Delta t=10^{-4}$~s) advancing the whole parameter grid simultaneously, as in the JR
sweeps. The maps report the
alpha-band ($8$--$12$~Hz) power of the $P_1$ and $P_2$ membrane potentials as the mean
squared magnitude of the Hann-windowed fast Fourier transform (FFT) over the band, in arbitrary units, over a
post-transient window. This measure includes autonomous and driven activity and
does not isolate a phase-coherent response to stimulation.

Two sweeps probe the column: an amplitude--beat sweep over
$(\varepsilon,\Delta f)$ at fixed carrier, and a carrier--beat sweep
over $(f_c,\Delta f)$ at fixed $\varepsilon$, built on the LaNMM module of
\cite{ruffini2025cfc}. The maps
(Figs.~\ref{fig:lanmm_p2},~\ref{fig:lanmm_p1}) show the \emph{signed} stimulation-induced
change, alpha power minus the unstimulated ($\varepsilon=0$) baseline, on a zero-centered
scale; over most of the $P_1$-driven grids the change is negative. For the
alpha resonance we instead drive the deep $P_1$ loop directly and sweep its input against
$\Delta f$, reading the lock-in (\ifsupp main-text\fi\ Eq.~\eqref{eq:lockin}) of $v_{P_1}$ at $\Delta f$.

The additional carrier comparison below uses the lock-in
of $v_{P_1}$, averaged at each carrier over 17 beats: $7$--$9$ and
$12.25$--$14$~Hz in $0.25$~Hz steps. These lie more than $1.5$~Hz from the
field-free alpha peak near $10.6$~Hz. We denote this arithmetic mean by $S(f_c)$.
The carrier grid is $25$--$80$~Hz in $1$~Hz steps, with equal tone amplitudes
$\varepsilon=1.5$~mV for $P_2$ drive and $0.5$~mV for $P_1$ drive. After $2$~s
of settling, we compare $3$ and $10$~s measurement windows. Field-free controls
give mean background amplitudes of $0.0048$ and $0.0012$~mV over the same beats.
This background includes a persistent component near $8$~Hz, so it is not solely
finite-window leakage. The statistic compares carrier profiles within each drive
condition; the unequal amplitudes do not support a direct sensitivity comparison
between targets.

\subsection{Frequency selection and carrier dependence}\label{sec:lanmm-results}
Driving either population with the two-tone field of Eq.~\eqref{eq:lanmm_drive}
and reading alpha-band power in the deep ($P_1$, alpha) and superficial ($P_2$,
gamma) pyramidal potentials recovers the beat resonance
(\S\ref{sec:lanmm-methods}, Fig.~\ref{fig:lanmm_setup}). Direct $P_1$ drive reproduces the JR
alpha-band resonance on both stable-focus and limit-cycle sides; that sweep
and its code are in the repository. The additional result is carrier selectivity
through the coupled gamma circuit.

Sweeping the field amplitude $\varepsilon$ and beat frequency $\Delta f$ at a fixed
carrier produces a region of enhanced alpha-band power: the response in $P_1$ is largest
when $\Delta f$ matches the deep alpha ($\sim$10~Hz) and grows with $\varepsilon$
(Figs.~\ref{fig:lanmm_p2}a,b and~\ref{fig:lanmm_p1}a,b). The dominant-frequency
contours identify where the output is close to the beat, without establishing
phase locking of the autonomous rhythm. Away from the enhancement region the drive
\emph{reduces} alpha-band power: with the field on $P_1$, $95\%$ of the amplitude--beat
grid and $97\%$ of the carrier--beat grid lie below the field-free baseline
(Fig.~\ref{fig:lanmm_p1}), the drive pulling the autonomous alpha out of the $8$--$12$~Hz
band or suppressing it; with the field on $P_2$ the fractions are $44\%$ and $31\%$
(Fig.~\ref{fig:lanmm_p2}).

The carrier dependence illustrates the two coupling limits of
\S\ref{sec:coupling}--\S\ref{sec:khz} within this model's $25$--$80$~Hz carrier
range. When the field is delivered to the superficial loop ($P_2$), alpha-band
power enhancement in $P_1$ is largest in the gamma carrier band with
$\Delta f$ near the $P_1$ alpha (Fig.~\ref{fig:lanmm_p2}c,d). At
$\Delta f=10$~Hz, the carrier maximum is $44$~Hz. This selectivity is consistent with a double
resonance: the fast loop acts as a tuned
radio-frequency (RF) front end whose dynamics affect the carrier dependence of
the deep response. With direct $P_1$ drive, a ridge near the column's own alpha
persists across the carrier range (Fig.~\ref{fig:lanmm_p1}c). That ridge contains
autonomous activity and is not a measurement of the demodulated component alone.

The lock-in also contains the autonomous rhythm when the beat is close to alpha;
its maximum over beats cannot isolate the driven response. We therefore compare
the off-alpha statistic $S(f_c)$ defined in \S\ref{sec:lanmm-methods}
(Table~\ref{tab:lanmm-carrier}).
With $P_2$ drive it peaks near $39$~Hz and falls to about $16$--$17\%$ of that
peak at a $25$~Hz carrier. With $P_1$ drive it peaks at $25$~Hz and spans about
$42$--$68\%$ of that peak over $30$--$80$~Hz. These are properties of the
specified amplitudes and beat average, not universal carrier-rejection factors.
Single-beat profiles have deeper minima, including a dip near a $45$~Hz carrier
under direct $P_1$ drive: the gamma circuit can affect the response in either
stimulation condition.

\begin{table}[t!]
\centering\small
\caption{\textbf{LaNMM carrier profiles measured away from alpha.}
$S(f_c)$ is the mean lock-in amplitude of $v_{P_1}$ over the 17 beats specified
in \S\ref{sec:lanmm-methods}. Percentages are normalized to each row's own carrier maximum;
the range and median use carriers from $30$ to $80$~Hz.}
\label{tab:lanmm-carrier}
\begin{tabular}{@{}llrrrr@{}}
\toprule
Drive & Window & Peak $S$ (mV) & Peak $f_c$ (Hz) & Range (\%) & Median (\%)\\
\midrule
$P_2$, $1.5$ mV & $3$ s  & $0.450$ & $39$ & $23$--$100$ & $42$\\
$P_2$, $1.5$ mV & $10$ s & $0.450$ & $39$ & $23$--$100$ & $42$\\
$P_1$, $0.5$ mV & $3$ s  & $0.092$ & $25$ & $42$--$68$  & $53$\\
$P_1$, $0.5$ mV & $10$ s & $0.091$ & $25$ & $42$--$67$  & $52$\\
\bottomrule
\end{tabular}
\end{table}

The carrier peaks and aggregate ranges persist with the longer window, although
individual normalized values change by up to $3.1$ percentage points for $P_2$
drive and $9.6$ for $P_1$ drive. The smallest $S$ exceeds its mean field-free
background by factors of about $16$ and $8$, respectively, with the $3$~s window,
and $62$ and $32$ with $10$~s. Some individual beat--carrier points remain near or below
background; the comparison concerns the mean over beats, not each grid point.
TI's kHz carriers lie outside this model's range. Within it, the carrier can
excite a network mode and affect the detected response, qualifying the inert-carrier
limit of \S\ref{sec:carrier} and linking it to frequency-tagging and
intermodulation paradigms \cite{gordon2019}.

\setcounter{figure}{0}
\begin{figure}[t!]
\centering
\includegraphics[width=0.86\textwidth]{fig_lanmm_field_p2_twotone.pdf}
\caption{\textbf{LaNMM frequency selection: two-tone TI field on the superficial $P_2$ (gamma)
loop.} The field enters the $P_2$ sigmoid argument as a membrane polarization $\varepsilon$
(mV) with all external drives flat. Color shows the signed change of alpha-band
($8$--$12$~Hz) power from the unstimulated $\varepsilon=0$ baseline, on a zero-centered
scale (arbitrary units). \emph{(a,b)} $P_1$ and $P_2$ enhancement vs.\
$\varepsilon$ and $\Delta f$ (carrier $40$~Hz), peaking at the $P_1$ alpha ($10.0$~Hz,
dashed). The black contour in \emph{(a)} is the frequency-dominance set
$|f_{\rm out}-\Delta f|<0.3$~Hz, drawn because band power
alone cannot separate a driven from an autonomous rhythm; it is a dominance screen, not a
phase-capture test. The $\varepsilon=0$ row is excluded, since with no drive the
autonomous $10.67$~Hz alpha already lies within the tolerance of two grid points (a
resolution floor of $0.5$~Hz); the set widens to $2.75$~Hz at $3$~mV. \emph{(c,d)} Enhancement vs.\ carrier and $\Delta f$ at
$\varepsilon=1.5$~mV: $P_1$ alpha-band enhancement is largest in the
gamma carrier band with $\Delta f$ near the $P_1$ alpha. At $\Delta f=10$~Hz,
the carrier maximum is $44$~Hz. The field-free
$P_2$ gamma peak is near $39.7$~Hz; the finite-drive maximum need not coincide with it.}
\label{fig:lanmm_p2}
\end{figure}



Within the $25$--$80$~Hz carrier range and at the finite polarizations used here,
the LaNMM shows alpha-band power changes concentrated near the intrinsic alpha,
amplitude-dependent widening of the frequency-dominance region, and enhanced
responses when superficial drive passes through the gamma circuit. The power
maps include autonomous and driven activity; they do not isolate TI efficacy.
The gamma-carrier preference, around $40$~Hz rather than kHz, remains a separate
prediction to test. Esmaeilpour et al.\ \cite{esmaeilpour2021} swept
$0.1$--$2$~kHz carriers at a fixed $5$~Hz beat against carbachol-induced gamma,
finding decreasing sensitivity with carrier frequency. Those carriers did not
sample the gamma band. Violante et al.\ \cite{violante2023} also used a fixed
$5$~Hz beat. Neither study tests this low-carrier resonance prediction.

\setcounter{figure}{1}
\begin{figure}[t!]
\centering
\includegraphics[width=0.85\textwidth]{fig_lanmm_setup.png}
\caption{\textbf{Stimulation setup in a LaNMM column.} The column couples a
superficial PING subcircuit (pyramidal $P_2$ + PV; gamma) and a deep JR
subcircuit (pyramidal $P_1$ + SS + SST; alpha). \emph{(a)} the TI field drives
$P_2$ while $P_1$ receives a flat drive; \emph{(b)} the TI field drives $P_1$
while $P_2$ is flat. Through the sigmoid nonlinearities and $P_1$--$P_2$ coupling,
the generated beat-frequency component is transferred to the deep loop: $P_2$ carries the
modulated fast carrier, while $P_1$ exposes the recovered low-frequency component (right
traces).}
\label{fig:lanmm_setup}
\end{figure}

\setcounter{figure}{6}
\begin{figure}[t!]
\centering
\includegraphics[width=0.86\textwidth]{fig_lanmm_field_p1_twotone.pdf}
\caption{\textbf{LaNMM frequency selection: two-tone TI field delivered directly to the deep
$P_1$ (alpha) loop.} As Fig.~\ref{fig:lanmm_p2} but with the field on $P_1$. The deep loop
shows frequency dominance at smaller fields, so the amplitude range is
$\varepsilon\le1$~mV here against $3$~mV there. These finite-amplitude maps extend
beyond the perturbative regime of the theory. Color: signed change of alpha-band power from the $\varepsilon=0$ baseline,
zero-centered; most of the grid is suppression. \emph{(a,b)} $P_1$ and $P_2$ enhancement vs.\ $\varepsilon$ and $\Delta f$; the
black contour in \emph{(a)} outlines the frequency-dominance set
$|f_{\rm out}-\Delta f|<0.3$~Hz (\S\ref{sec:vocab}), computed from the dominant
output frequency because band power alone cannot separate a driven from an autonomous
rhythm; the $\varepsilon=0$ row is excluded (resolution floor $0.5$~Hz). The set widens to
$1.25$~Hz at $\varepsilon=1$~mV.
\emph{(c,d)} Enhancement vs.\ carrier frequency and $\Delta f$ at $\varepsilon=0.5$~mV.
The ridge near alpha persists across the carrier band and includes the autonomous
rhythm. A separate lock-in comparison away from alpha, including its window and
background controls, is given in Table~\ref{tab:lanmm-carrier}; these panels show
band-power changes.}
\label{fig:lanmm_p1}
\end{figure}

